\documentclass[12pt]{article}
\usepackage{notes-project}
\usepackage{braket}

\begin{document}

\begin{center}
  {\LARGE\bfseries Chapter 6: Symmetries}\\[0.25em]
  {\large Quantum Mechanics}\\[0.15em]
\end{center}

\section{Introduction to Symmetry}

\begin{dfn}[Symmetry]{symmetry}
  A \emph{symmetry} of a system is some transformation, such that the system is invariant before and after the transformation.
\end{dfn}

\begin{dfn}[Symmetry of Quantum System]{symmetry-qm}
  A \emph{symmetry} of a quantum system is a transformation that leaves the Hamiltonian invariant.
\end{dfn}

\begin{dfn}[Types of Symmetry]{symmetry-types}
  There are several types of symmetries in quantum mechanics, including:
  \begin{itemize}[nosep]
    \item \emph{Translation symmetry}: The system is invariant under some spatial translations.
    \item \emph{Rotation symmetry}: The system is invariant under some rotations.
    \item \emph{Reflection symmetry}: The system is invariant under reflections on some axis/planes.
    \item \emph{Inversion symmetry}: The system is invariant under inversion of all coordinates.
    \item \emph{Time-reversal symmetry}: The system is invariant under the reversal of time.
  \end{itemize}
\end{dfn}

\begin{thm}[Symmetry Condition on the Hamiltonian]{symmetry-condition-hamiltonian}
  An operator $\hat{U}$ realizes a \xref{def:symmetry-qm} of a system with Hamiltonian $\hat{H}$ if and only if $\hat{H}$ is invariant under \xref{ch03::def:conjugation-action}:
  \[
    \hat{U}^{\dagger}\hat{H}\hat{U} = \hat{H}, \qquad \text{equivalently (for linear $\hat{U}$)} \qquad [\hat{U}, \hat{H}] = 0.
  \]
  For antiunitary $\hat{U}$ (e.g.\ time reversal), the second form must be replaced by the appropriate antilinear condition.
\end{thm}


\section{Translation Operator}

\begin{dfn}[Translation Operator]{translation-operator}
  Intuitively, a \emph{translation operator} $\hat{T}(\vec{r}_0)$ (translates a state by a vector $\vec{r}_0$) acts on a state $\ket{\psi}$ should let $\ket{\psi}$ be translated by a vector $\vec{r}$, i.e.
  \[\hat{T}(\vec{r}_0)\ket{\psi} = \ket{\psi(\vec{r} - \vec{r}_0)}\]
\end{dfn}

\begin{dfn}[Translation Operator (1D)]{translation-operator-1d}
  In one dimension, the translation operator $\hat{T}(a)$ shifts a state by a real distance $a \in \mathbb{R}$. Its action on a position-space wavefunction is
  \[
    \hat{T}(a)\,\psi(x) = \psi(x - a).
  \]
\end{dfn}

\begin{thm}[Exponential Form of the Translation Operator]{translation-exponential}
  The one-dimensional translation operator can be written as:
  \[\hat{T}(a) = e^{-i a \hat{p}/\hbar}\]
\end{thm}

\begin{proof}
  Expand the exponential in a \xref[Taylor series]{ch03::def:operator-function} at $0$ 
  \begin{align*}
  \hat{T}(a) &= e^{-i a \hat{p}/\hbar} = \sum_{n=0}^{\infty} \frac{(-ia/\hbar)^n}{n!} \cdot \hat{p}^n = \sum_{n=0}^{\infty} \frac{(-i a \hat{p})^n}{n! \hbar^{n}}
  \end{align*} 

  According to the \xref[definition of the momentum operator]{ch02::def:momentum-operator-position}, $\hat{p} = -i\hbar \frac{d}{dx}$, so we can substitute this into the series expansion:
  \begin{align*}
  \hat{T}(a) &= \sum_{n=0}^{\infty} \frac{(-i a (-i\hbar \frac{d}{dx}))^n}{n! \hbar^{n}} = \sum_{n=0}^{\infty} \frac{(- a ( \frac{d}{dx}))^n}{n!} 
  \end{align*}

  So, 
  \[\hat{T}(a) \psi(x) = \sum_{n=0}^{\infty} \frac{(- a ( \frac{d}{dx}))^n}{n!} \psi(x) = \sum_{n=0}^{\infty} \frac{(- a)^n}{n!} \psi^{(n)}(x)\] 
  
  Then, we do a Taylor expansion of $\psi(x)$ centered at $x + a$: 
  \[\psi(x) = \sum_{n=0}^{\infty} \frac{\psi^{(n)}(x + a)}{n!} (x - (x + a))^n = \sum_{n=0}^{\infty} \frac{\psi^{(n)}(x + a)}{n!} (-a)^n = \sum_{n=0}^{\infty} \frac{(-a)^n}{n!} \psi^{(n)}(x + a) \] 

  Substituting $x \to x - a$ in the previous identity,
  \[\psi(x - a) = \sum_{n=0}^{\infty} \frac{(-a)^n}{n!} \psi^{(n)}((x - a) + a) = \sum_{n=0}^{\infty} \frac{(-a)^n}{n!} \psi^{(n)}(x).\]

  Which means, 
  \[\hat{T}(a) \psi(x) = \psi(x - a)\] 
\end{proof}

\begin{thm}[$\hat T(a)$ is a Unitary Operator]{translation-unitary}
  The translation operator $\hat{T}(a)$ is a \xref[unitary operator]{ch03::def:unitary-operator}, i.e. $\hat{T}^{\dagger}(a) = \hat{T}^{-1}(a) = \hat{T}(-a)$.
\end{thm}

\begin{rem}[Intuition about ]

\end{rem}

\begin{proof}
  {\color{blue}\textbf{TODO}: Proof to be written after clarifying preferred approach with the lecturer --- direct from the action definition (\xref{def:translation-operator-1d}), or via the exponential form (\xref{thm:translation-exponential}). Email sent \today.}
\end{proof}

\end{document}
